%------------------------------------------------------------------------------%
\documentclass[fleqn,utf8,aspectratio=169,12pt]{beamer}

\usepackage{gaal}

\title{Autovalores e Autovetores}

%------------------------------------------------------------------------------%
\begin{document}

\addtitle

%------------------------------------------------------------------------------%
\section{Transformações Lineares}

%------------------------------------------------------------------------------%
\begin{frame}{Transformações Lineares}
  Uma \emph{transformação linear} é uma função
  \[
    \pause
    T\colon \R^n \to \R^m
  \]
  \pause
  que satisfaz as seguintes propriedades
  \[
    \pause
    T(\vec{u}+\vec{v}) = T(\vec{u}) + T(\vec{v})
  \]
  \[
    \pause
    T(\alpha\vec{v})=\alpha T(\vec{v})
  \]
  \pause
  para quaisquer $\vec{u}$, $\vec{v}\in\R^n$ e $\alpha\in\R$
\end{frame}

%------------------------------------------------------------------------------%
\begin{frame}{Representação Matricial}
  Toda \emph{transformação linear}
  \[
    T\colon \R^n \to \R^m
  \]

  \pause
  \medskip
  pode ser representada por uma \emph{matriz} $A$
  \pause
  de ordem $m\times n$
  \pause
  \[
    T(\vec{x}) = A\vec{x}
  \]
\end{frame}

%------------------------------------------------------------------------------%
\begin{frame}{Representação Matricial}
  Escrevendo
  \[
    A
    \pause
    = \begin{bmatrix}[c]
      a_{11} & \cdots & a_{1n} \\
      \vdots & \ddots & \vdots \\
      a_{m1} & \cdots & a_{mn}
    \end{bmatrix}
    \qquad\qquad\pause
    \vec{x} =
    \begin{bmatrix}[c]
      x_1    \\
      \vdots \\
      x_n
    \end{bmatrix}
  \]

  \pause
  \medskip
  A transformação linear fica
  \[
    \pause
    T(\vec{x}) =
    \begin{bmatrix}[c]
      a_{11}x_1+\cdots+a_{1n}x_n \\
      \vdots                     \\
      a_{m1}x_1+\cdots+a_{mn}x_n
    \end{bmatrix}
  \]
\end{frame}

%------------------------------------------------------------------------------%
%------------------------------------------------------------------------------%
\begin{question}
  Represente matricialmente a transformação linear
  \[
    T\colon \R^2 \to \R^2
  \]
  definida por
  \[
    T(x, y) = ( 2x, y )
  \]

  \medskip
  Avalie $T$ em
  \(
    \vec{v} = \vetor{1}{2}
  \)
\end{question}

%------------------------------------------------------------------------------%
\begin{resolution}{Representação Matricial}
    A representação matricial da transformação linear
    \[
        T(x, y) = (2x, y)
    \]
    \pause
    é
    \[
      T \vetor{x}{y}
      \pause
      = \begin{bmatrix}
          2 & 0\\
          0 & 1
      \end{bmatrix}
      \vetor{x}{y}
    \]
\end{resolution}

%------------------------------------------------------------------------------%
\begin{resolution}{Avalianção}
  Avaliando $T$ no vetor
  \[
    \vec{v} = \vetor{1}{2}
  \]
  \pause
  obtemos
  \[
    T(\vec{v})
    \pause
    = \begin{bmatrix}
        2 & 0\\
        0 & 1
    \end{bmatrix}
    \vetor{1}{2}
    \pause
    = \vetor{2}{2}
  \]
\end{resolution}

%------------------------------------------------------------------------------%
\begin{resolution}{Interpretação Geométrica}
  A transformação
  \[
      T(x, y) = (2x, y)
  \]
  \pause
  duplica a coordenada horizontal e mantém a coordenada vertical
\end{resolution}

%------------------------------------------------------------------------------%
\begin{resolution}{Interpretação Geométrica}
  O vetor é transformado em outro vetor

\pause
\begin{center}
  \begin{tikzpicture}[scale=1.5,>=stealth]

    \draw[->] (-0.5,0)--(4.2,0) node[right] {$x$};
    \draw[->] (0,-0.5)--(0,3.2) node[above] {$y$};

    \pause
    \draw[->,thick,blue] (0,0)--(1,2) node[above left ] {$\vec{v}$};
    \draw[dashed] (1,0)--(1,2);

    \pause
    \draw[dashed] (2,0)--(2,2);
    \draw[->,thick,red]  (0,0)--(2,2) node[above right] {$T(\vec{v})$};

  \end{tikzpicture}
\end{center}
\end{resolution}

%------------------------------------------------------------------------------%
\endinput

%------------------------------------------------------------------------------%
\begin{question}
  Represente matricialmente a transformação linear
  \[
    T\colon \R^2\longrightarrow\R^2
  \]
  definida por
  \[
    T(x, y) = (3x, 2y)
  \]

  \medskip
  Avalie $T$ em
  \(
    \vec{u} = \vetor{1}{1}
  \),
  \(
    \vec{v} = \vetor{1}{0}
  \)
  e
  \(
    \vec{v} = \vetor{0}{1}
  \)
\end{question}

%------------------------------------------------------------------------------%
\begin{resolution}{Representação Matricial}
    A representação matricial da transformação linear
    \[
      T(x, y) = (3x, 2y)
    \]
    \pause
    é
    \[
      T \vetor{x}{y} =
      \begin{bmatrix}
        3 & 0 \\
        0 & 2
      \end{bmatrix}
      \vetor{x}{y}
    \]
\end{resolution}

%------------------------------------------------------------------------------%
\begin{resolution}{Avalianção do Vetor $\vec{u}$}
  Avaliando $T$ no vetor
  \[
    \vec{u} = \vetor{1}{1}
  \]
  \pause
  obtemos
  \[
    T(\vec{u})
    \pause
    = \begin{bmatrix}
        3 & 0\\
        0 & 2
    \end{bmatrix} \vetor{1}{1}
    \pause
    = \vetor{3}{2}
  \]
\end{resolution}

%------------------------------------------------------------------------------%
\begin{resolution}{Avalianção do Vetor $\vec{v}$}
  Avaliando $T$ no vetor
  \[
    \vec{v} = \vetor{1}{0}
  \]
  \pause
  obtemos
  \[
    T(\vec{v})
    \pause
    = \begin{bmatrix}
        3 & 0\\
        0 & 2
    \end{bmatrix} \vetor{1}{0}
    \pause
    = \vetor{3}{0}
    \pause
    = 3 \vec{v}
  \]

  \pause
  \emph{A direção do vetor foi preservada}
\end{resolution}

%------------------------------------------------------------------------------%
\begin{resolution}{Avalianção do Vetor $\vec{w}$}
  Avaliando $T$ no vetor
  \[
    \vec{w} = \vetor{0}{1}
  \]
  \pause
  obtemos
  \[
    T(\vec{w})
    \pause
    = \begin{bmatrix}
        3 & 0\\
        0 & 2
    \end{bmatrix} \vetor{0}{1}
    \pause
    = \vetor{0}{2}
    \pause
    = 2 \vec{w}
  \]

  \pause
  \emph{A direção do vetor foi preservada}
\end{resolution}

%------------------------------------------------------------------------------%
\endinput


%------------------------------------------------------------------------------%
\section{Autovetores}

%------------------------------------------------------------------------------%
\begin{frame}{Autovetores de uma Transformação Linear}
  Em uma transformação linear

  \pause
  alguns vetores especiais têm sua direção preservada

  \pause
  Para esses vetores, a imagem é um múltiplo escalar do próprio vetor
  \[
    \pause
    T(\vec{v}) = \lambda\vec{v}
  \]

  \pause
  Representação Matricial

  \pause
  Como \(T(\vec{v}) = A\vec{v}\),
  \pause
  podemos escrever
  \[
    A\vec{v} = \lambda\vec{v}
  \]
\end{frame}

%------------------------------------------------------------------------------%
\begin{frame}{Autovalores}
  Um \emph{autovetor} é um vetor cuja direção é preservada

  \pause
  quando o multiplicamos por uma matriz
\end{frame}

%------------------------------------------------------------------------------%
\begin{frame}{Definição}
  Dada uma matriz \(A\in\R^{n\times n}\)

  \pause
  \bigskip
  Um vetor \(\vec{v}\neq\vec{0}\) é um \emph{autovetor} de \(A\) quando
  \[
    \pause
    A\vec{v} = \lambda\vec{v}
  \]

  \pause
  \medskip
  \(\lambda\) é o \emph{autovalor} de \(A\) associado a \(\vec{v}\)
\end{frame}

%------------------------------------------------------------------------------%
\begin{frame}{Equação Característica}
  Temos $n$ equações

  \pause
  \bigskip
  Precisamos determinar $n+1$ incognitas

  \pause
  \qquad As $n$ componentes do autovetor $\vec{v}$

  \pause
  \qquad O autovalor $\lambda$

  \pause
  \bigskip
  Precisamos de mais uma equação
\end{frame}

%------------------------------------------------------------------------------%
\begin{frame}{Equação Característica}
\begin{align*}
  \onslide<+->{A\vec{v}                  & = \lambda\vec{v} \\}
  \onslide<+->{A\vec{v} - \lambda\vec{v} & = \vec{0}        \\}
  \onslide<+->{(A-\lambda I)\vec{v}      & = \vec{0}}
\end{align*}

  \medskip
  \onslide<+->{Para existir uma solução não nula}
  \[
    \onslide<+->{\det(A-\lambda I) = 0}
  \]

  \smallskip
  \onslide<+->{Essa é a \emph{equação característica} de \(A\)}
\end{frame}

%------------------------------------------------------------------------------%
\begin{frame}{Polinômio Característico}
  O \emph{polinômio característico} é
  \[
    \pause
    p_{A}(\lambda)
    \pause
    = \det(A-\lambda I)
    \pause
    = a_n\lambda^n
    + a_{n-1}\lambda^{n-1}
    + \cdots
    + a_1 \lambda
    + a_0
  \]

  \pause
  Suas raízes são os \emph{autovalores} de \(A\)
\end{frame}

%------------------------------------------------------------------------------%
%------------------------------------------------------------------------------%
\begin{question}
  Determine os autovalores da matriz
  \(
    A =
    \begin{bmatrix}
      4 & 1 \\
      2 & 3
    \end{bmatrix}
  \)
\end{question}

%------------------------------------------------------------------------------%
\begin{resolution}{Polinômio Característica}
  \begin{align*}
    \onslide<+->{p(\lambda)}
    \onslide<+->{& = \det(A - \lambda I)}
    \\
    \onslide<+->{& = \det\left(
          \begin{bmatrix}
            4 & 1 \\
            2 & 3
          \end{bmatrix}
          - \lambda
          \begin{bmatrix}
            1 & 0 \\
            0 & 1
          \end{bmatrix}
    \right)}
    \\
    \onslide<+->{& = \begin{vmatrix}[c]
          4-\lambda & 1         \\
          2         & 3-\lambda
        \end{vmatrix}}
    \\
    \onslide<+->{& = (4-\lambda)(3-\lambda) - 1 \times 2  \\}
    \onslide<+->{& = 12 - 4\lambda - 3\lambda + \lambda^2 - 2        \\}
    \onslide<+->{& = \lambda^2 - 7\lambda + 10}
  \end{align*}
\end{resolution}

%------------------------------------------------------------------------------%
\begin{resolution}{Equação Característica}
  \[
    p(\lambda)
    \;=\; \lambda^2 - 7\lambda + 10
    \pause
    \;=\; 0
  \]

  \[
    \pause
    \Delta
    \pause
    \;=\; b^2 - 4ac
    \pause
    \;=\; (-7)^2 - 4 \times 1 \times 10
    \pause
    \;=\; 49 - 40
    \pause
    \;=\; 9
  \]

  \[
    \pause
    \sqrt{\Delta}
    \pause
    \;=\; \sqrt{9}
    \pause
    \;=\; 3
  \]

  \[
    \pause
    \lambda_1
    \pause
    \;=\; \frac{-b + \sqrt{\Delta}}{2a}
    \pause
    \;=\; \frac{7 + 3}{2}
    \pause
    \;=\; 5
  \]

  \[
    \pause
    \lambda_2
    \pause
    \;=\; \frac{-b - \sqrt{\Delta}}{2a}
    \pause
    \;=\; \frac{7 - 3}{2}
    \pause
    \;=\; 2
  \]
\end{resolution}

%------------------------------------------------------------------------------%
\endinput

%------------------------------------------------------------------------------%
\begin{question}
  Determine os autovetores da matriz
  \(
    A =
    \begin{bmatrix}
      4 & 1 \\
      2 & 3
    \end{bmatrix}
  \)

  \bigskip
  Note que, no exercício anterior, determinamos que os autovalores de $A$ são
  \[
    \lambda_1 = 5
    \qquad
    \lambda_2 = 2
  \]
\end{question}

%------------------------------------------------------------------------------%
\begin{resolution}{Autovetores Associados a \(\lambda_1 = 5\)}
\begin{columns}
\column{0.5\textwidth}
  \begin{align*}
    \onslide<+->{(A-\lambda_1 I)\vec{v} & = \vec{0} \\}
    \onslide<+->{(A-5I)\vec{v} & = \vec{0}}
    \\
    \onslide<+->{\left(
      \begin{bmatrix}
        4 & 1 \\
        2 & 3
      \end{bmatrix}
      -
      \begin{bmatrix}
        5 & 0 \\
        0 & 5
      \end{bmatrix}
    \right)
        \begin{bmatrix} x \\ y \end{bmatrix}
    & = \begin{bmatrix} 0 \\ 0 \end{bmatrix}}
    \\
    \onslide<+->{\begin{bmatrix}
      -1  &  1  \\
       2  & -2
    \end{bmatrix}
        \begin{bmatrix} x \\ y \end{bmatrix}
    & = \begin{bmatrix} 0 \\ 0 \end{bmatrix}}
    \\
    \onslide<+->{\begin{cases}
      -x +  y = 0 \\
      2x - 2y = 0
    \end{cases}}
  \end{align*}
\column{0.5\textwidth}
  \begin{align*}
    \onslide<+->{-x + y & = 0 \\}
    \onslide<+->{     y & = x}
  \end{align*}
  \onslide<+->{Assim}
  \(
    \onslide<+->{\vec{v}}
    \onslide<+->{=   \begin{bmatrix} x \\ x \end{bmatrix}}
    \onslide<+->{= x \begin{bmatrix} 1 \\ 1 \end{bmatrix}}
  \)

  \medskip
  \onslide<+->{Os autovetores são \\
  os múltiplos não nulos}
  \onslide<+->{de
  \[
    \vec{v}_1 =
    \begin{bmatrix}
      1 \\
      1
    \end{bmatrix}
  \]}
\end{columns}
\end{resolution}


%------------------------------------------------------------------------------%
\begin{resolution}{Autovetores Associados a \(\lambda_2 = 2\)}
\begin{columns}
\column{0.5\textwidth}
  \begin{align*}
    \onslide<+->{(A-\lambda_2 I)\vec{v} & = \vec{0} \\}
    \onslide<+->{(A-2I)\vec{v} & = \vec{0} }
    \\
    \onslide<+->{\left(
      \begin{bmatrix}
        4 & 1 \\
        2 & 3
      \end{bmatrix}
      -
      \begin{bmatrix}
        2 & 0 \\
        0 & 2
      \end{bmatrix}
    \right)
        \begin{bmatrix} x \\ y \end{bmatrix}
    & = \begin{bmatrix} 0 \\ 0 \end{bmatrix}}
    \\
    \onslide<+->{\begin{bmatrix}
      2 & 1  \\
      2 & 1
    \end{bmatrix}
        \begin{bmatrix} x \\ y \end{bmatrix}
    & = \begin{bmatrix} 0 \\ 0 \end{bmatrix}}
    \\
    \onslide<+->{\begin{cases}
      2x + y = 0 \\
      2x + y = 0
    \end{cases}}
  \end{align*}
\column{0.5\textwidth}
  \[
    \onslide<+->{y = -2x}
  \]
  \onslide<+->{Assim}
  \(
    \onslide<+->{\vec{v}}
    \onslide<+->{= \begin{bmatrix}
      x \\
      -2x
    \end{bmatrix}}
    \onslide<+->{= x
    \begin{bmatrix}
      1 \\
      -2
    \end{bmatrix}}
  \)

  \medskip
  \onslide<+->{Os autovetores são \\
  os múltiplos não nulos}
  \onslide<+->{de
  \[
    \vec{v}_2 =
    \begin{bmatrix}
      1 \\
      -2
    \end{bmatrix}
  \]}
\end{columns}
\end{resolution}

%------------------------------------------------------------------------------%
\endinput



%------------------------------------------------------------------------------%
\begin{question}
  Determine os autovalores e autovetores da matriz
  \[
    A =
    \begin{bmatrix}
      0 & -1 \\
      1 & 0
    \end{bmatrix}
  \]
\end{question}

%------------------------------------------------------------------------------%
\begin{resolution}{Autovalores}
  \onslide<+->{Os autovalores são as raízes do polinômio característico}
  \[
    \onslide<+->{p(\lambda)}
    \onslide<+->{= \det(A-\lambda I)}
    \onslide<+->{= \begin{vmatrix}
        -\lambda & -1       \\
        1        & -\lambda
      \end{vmatrix}}
    \onslide<+->{= \lambda^{2}+1}
  \]
  \begin{align*}
    \onslide<+->{\lambda^{2}+1 & =  0 \\}
    \onslide<+->{\lambda^{2}   & = -1}
  \end{align*}
  \[
    \onslide<+->{\nexists \, \lambda \in \R}
  \]
\end{resolution}

%------------------------------------------------------------------------------%
\begin{resolution}{Interpretação geométrica}
  Matriz de rotação de $90^\circ$ no sentido anti-horário

\begin{center}
  \begin{tikzpicture}[scale=2,>=stealth]
    \draw[->] (-1.8,0)--(1.8,0) node[right] {$x$};
    \draw[->] (0,-0.5)--(0,1.8) node[above] {$y$};

    \pause
    \draw[->,thick,blue] (0,0)--(1.2,0.5) node[right] {$\vec{v}$};

    \pause
    \draw[->,thick,red] (0,0)--(-0.5,1.2) node[above left] {$A\vec{v}$};

    \draw[->,dashed] (0.65,0.27) arc[start angle=22.6,end angle=112.6,radius=0.7];
  \end{tikzpicture}
\end{center}
\end{resolution}

%------------------------------------------------------------------------------%
\endinput


%------------------------------------------------------------------------------%
\section{Autoespaços}

%------------------------------------------------------------------------------%
\begin{frame}{Autoespaços}
  O \emph{autoespaço} associado a um autovalor \(\lambda\)

  \pause
  contém o \emph{vetor nulo} e todos os autovetores associados a \(\lambda\)

  \pause
  \[
    E(\lambda) =
    \left\{\,
      \vec{v} \in \R^n \st
      (A-\lambda I)\vec{v} = \vec{0}
    \,\right\}
  \]
\end{frame}

%------------------------------------------------------------------------------%
%------------------------------------------------------------------------------%
\begin{question}
  Determine os autovalores, os autovetores e os autoespaços da matriz
  \[
    B =
    \begin{bmatrix}
      2 & 1 & 0 \\
      0 & 2 & 0 \\
      0 & 0 & 5
    \end{bmatrix}
  \]
\end{question}

%------------------------------------------------------------------------------%
\begin{resolution}{Autovalores}
  \onslide<+->{Polinômio característico}
  \begin{align*}
    \onslide<+->{p(\lambda)}
    \onslide<+->{& = \det(B - \lambda I) \\}
    \onslide<+->{& = \det\left(
          \begin{bmatrix}
            2 & 1 & 0 \\
            0 & 2 & 0 \\
            0 & 0 & 5
          \end{bmatrix}
          -
          \begin{bmatrix}
            \lambda & 0       & 0       \\
            0       & \lambda & 0       \\
            0       & 0       & \lambda
          \end{bmatrix}
          \right)}
    \onslide<+->{\;=\; \begin{vmatrix}
      2-\lambda & 1         & 0         \\
      0         & 2-\lambda & 0         \\
      0         & 0         & 5-\lambda
    \end{vmatrix}}
    \\
    \onslide<+->{& = (2-\lambda)^2(5-\lambda)}
  \end{align*}
  \onslide<+->{Os autovalores são}
  \[
    \onslide<+->{\lambda_{1} = 2
    \quad\text{ e }\quad
    \lambda_{2} = 5}
  \]
\end{resolution}

%------------------------------------------------------------------------------%
\begin{resolution}{Autovetores de \(\lambda=2\)}
\begin{columns}
\column{0.5\textwidth}
  \begin{align*}
    \onslide<+->{(B-2I)\vec{v} & = \vec{0}}
    \\[1mm]
    \onslide<+->{\begin{bmatrix}[c]
      2-2 & 1   & 0   \\
      0   & 2-2 & 0   \\
      0   & 0   & 5-2
    \end{bmatrix}
    \begin{bmatrix} x \\ y \\ z \end{bmatrix}
    & =
    \begin{bmatrix} 0 \\ 0 \\ 0 \end{bmatrix}}
    \\[3mm]
    \onslide<+->{\begin{bmatrix}
      0 & 1 & 0 \\
      0 & 0 & 0 \\
      0 & 0 & 3
    \end{bmatrix}
    \begin{bmatrix} x \\ y \\ z \end{bmatrix}
    & = \begin{bmatrix} 0 \\ 0 \\ 0 \end{bmatrix}}
   \end{align*}
\column{0.4\textwidth}
  \begin{align*}
    \onslide<+->{x & = t \in \R \\}
    \onslide<+->{y & = 0 \\}
    \onslide<+->{z & = 0}
  \end{align*}
  \onslide<+->{Autovetor associado a $\lambda = 2$}
  \[
    \onslide<+->{\vec{v}_1
    = \vetortri{1}{0}{0}}
  \]
\end{columns}
\end{resolution}

%------------------------------------------------------------------------------%
\begin{resolution}{Autoespaço de \(\lambda=2\)}
  \[
    E(2)
    \pause
    = \left\{\vetortri{t}{0}{0} \st t \in \R \right\}
    \pause
    = \ger\left\{ \vetortri{1}{0}{0} \right\}
  \]
\end{resolution}

%------------------------------------------------------------------------------%
\begin{resolution}{Autovetores de \(\lambda=5\)}
\begin{columns}
\column{0.5\textwidth}
  \begin{align*}
    \onslide<+->{(B-5I)\vec{v} & = \vec{0}}
    \\[1mm]
    \onslide<+->{\begin{bmatrix}[c]
      2-5 & 1   & 0   \\
      0   & 2-5 & 0   \\
      0   & 0   & 5-5
    \end{bmatrix}
    \begin{bmatrix} x \\ y \\ z \end{bmatrix}
    & = \begin{bmatrix} 0 \\ 0 \\ 0 \end{bmatrix}}
    \\[3mm]
    \onslide<+->{\begin{bmatrix}
      -3 & 1  & 0 \\
      0  & -3 & 0 \\
      0  & 0  & 0
    \end{bmatrix}
    \begin{bmatrix} x \\ y \\ z \end{bmatrix}
    & = \begin{bmatrix} 0 \\ 0 \\ 0 \end{bmatrix}}
   \end{align*}
\column{0.4\textwidth}
  \begin{align*}
    \onslide<+->{-3y   & = 0     \\}
    \onslide<+->{-3x+y & = 0     \\}
    \onslide<+->{x     & = y = 0 \\}
    \onslide<+->{z     & = t \in \R}
  \end{align*}
  \onslide<+->{Autovetor associado a $\lambda = 5$}
  \[
    \onslide<+->{\vec{v}_2
    = \vetortri{0}{0}{1}}
  \]
\end{columns}
\end{resolution}

%------------------------------------------------------------------------------%
\begin{resolution}{Autoespaço de \(\lambda=5\)}
  \[
    E(5)
    \pause
    = \left\{\vetortri{0}{0}{t} \st t \in \R \right\}
    \pause
    = \ger\left\{ \vetortri{0}{0}{1} \right\}
  \]
\end{resolution}

%------------------------------------------------------------------------------%
\begin{resolution}{Resposta}
  Autovalores
  \(
    \lambda_{1} = 2
    \quad\text{ e }\quad
    \lambda_{2} = 5
  \)

  \pause
  Autovetor e autoespaço associados a \(\lambda_{1} = 2\)
  \[
    \vec{v}_1 = \vetortri{1}{0}{0}
    \quad\text{ e }\quad
    E(2) = \ger\left\{ \vetortri{1}{0}{0} \right\}
  \]

  \pause
  Autovetor e autoespaço associados a \(\lambda_{2} = 5\)
  \[
    \vec{v}_2 = \vetortri{0}{0}{1}
    \quad\text{ e }\quad
    E(5) = \ger\left\{ \vetortri{0}{0}{1} \right\}
  \]
\end{resolution}

%------------------------------------------------------------------------------%
\endinput

%------------------------------------------------------------------------------%
\begin{question}
  Determine os autovalores, os autovetores e os autoespaços da matriz
  \[
    A =
    \begin{bmatrix}
      2 & 1 & 0  \\
      0 & 3 & 0  \\
      0 & 0 & -1
    \end{bmatrix}
  \]
\end{question}

%------------------------------------------------------------------------------%
\begin{resolution}{Autovalores}
  \onslide<+->{Os autovalores são as raízes do polinômio característico}
  \begin{align*}
      \onslide<+->{p(\lambda)}
      \onslide<+->{& = \det(A-\lambda I)}
      \\
      \onslide<+->{& = \det\left( \begin{bmatrix}
          2 & 1 & 0  \\
          0 & 3 & 0  \\
          0 & 0 & -1
        \end{bmatrix}
      - \lambda\begin{bmatrix}
          1 & 0 & 0  \\
          0 & 1 & 0  \\
          0 & 0 & 1
        \end{bmatrix}
      \right)}
      \onslide<+->{\;=\; \begin{vmatrix}[c]
        2-\lambda & 1         & 0          \\
        0         & 3-\lambda & 0          \\
        0         & 0         & -1-\lambda
      \end{vmatrix}}
      \\
      \onslide<+->{& = (2-\lambda)(3-\lambda)(-1-\lambda)}
  \end{align*}

  \onslide<+->{Portanto
  \[
    \lambda_{1} = 2
    \qquad
    \lambda_{2} = 3
    \qquad
    \lambda_{3} = -1
  \]}
\end{resolution}

%------------------------------------------------------------------------------%
\begin{resolution}{Autovetores para $\lambda=2$}
\begin{columns}
\column{0.7\textwidth}
  \begin{align*}
    \onslide<+->{(A-2I)\vec{x} & = \vec{0}}
    \\[3mm]
    \onslide<+->{\left(
      \begin{bmatrix}
        2 & 1 & 0  \\
        0 & 3 & 0  \\
        0 & 0 & -1
      \end{bmatrix}
    - \begin{bmatrix}
        2 & 0 & 0  \\
        0 & 2 & 0  \\
        0 & 0 & 2
      \end{bmatrix}
    \right)
    \vetortri{x}{y}{z}
    & = \vetortri{0}{0}{0}}
    \\[3mm]
    \onslide<+->{\begin{bmatrix}
      0 & 1 & 0  \\
      0 & 1 & 0  \\
      0 & 0 & -3
    \end{bmatrix}
    \vetortri{x}{y}{z}
    & = \vetortri{0}{0}{0}}
  \end{align*}

\column{0.3\textwidth}
  \begin{align*}
    \onslide<+->{x & \in \R \\}
    \onslide<+->{y & = 0    \\}
    \onslide<+->{z & = 0}
  \end{align*}

  \onslide<+->{Autovetor
  \[
    \vec{v}_1 = \vetortri{1}{0}{0}
  \]}
\end{columns}
\end{resolution}

%------------------------------------------------------------------------------%
\begin{resolution}{Autovetores para $\lambda=3$}
\begin{columns}
\column{0.7\textwidth}
  \begin{align*}
    \onslide<+->{(A-3I)\vec{x} & = \vec{0}}
    \\[3mm]
    \onslide<+->{\left(
      \begin{bmatrix}
        2 & 1 & 0  \\
        0 & 3 & 0  \\
        0 & 0 & -1
      \end{bmatrix}
    - \begin{bmatrix}
        3 & 0 & 0  \\
        0 & 3 & 0  \\
        0 & 0 & 3
      \end{bmatrix}
    \right)
    \vetortri{x}{y}{z}
    & = \vetortri{0}{0}{0}}
    \\[3mm]
    \onslide<+->{\begin{bmatrix}
      -1 & 1 & 0  \\
      0  & 0 & 0  \\
      0  & 0 & -4
    \end{bmatrix}
    \vetortri{x}{y}{z}
    & = \vetortri{0}{0}{0}}
  \end{align*}

\column{0.3\textwidth}
  \begin{align*}
    \onslide<+->{-x+y & = 0 \\}
    \onslide<+->{z    & = 0 \\}
    \onslide<+->{y    & = x}
  \end{align*}

  \onslide<+->{Autovetor
  \[
    \vec{v}_2 = \vetortri{1}{1}{0}
  \]}
\end{columns}
\end{resolution}

%------------------------------------------------------------------------------%
\begin{resolution}{Autovetores para $\lambda=-1$}
\begin{columns}
\column{0.7\textwidth}
  \begin{align*}
    \onslide<+->{(A-(-1)I)\vec{x} & = \vec{0}}
    \\[3mm]
    \onslide<+->{\left(
      \begin{bmatrix}
        2 & 1 & 0  \\
        0 & 3 & 0  \\
        0 & 0 & -1
      \end{bmatrix}
    + \begin{bmatrix}
        1 & 0 & 0  \\
        0 & 1 & 0  \\
        0 & 0 & 1
      \end{bmatrix}
    \right)
    \vetortri{x}{y}{z}
    & = \vetortri{0}{0}{0}}
    \\[3mm]
    \onslide<+->{\begin{bmatrix}
      3 & 1 & 0 \\
      0 & 4 & 0 \\
      0 & 0 & 0
    \end{bmatrix}
    \vetortri{x}{y}{z}
    & = \vetortri{0}{0}{0}}
  \end{align*}

\column{0.3\textwidth}
  \begin{align*}
    \onslide<+->{z    & \in \R \\}
    \onslide<+->{3x+y & = 0    \\}
    \onslide<+->{y    & = 0    \\}
    \onslide<+->{x    & = 0}
  \end{align*}

  \onslide<+->{Autovetor
  \[
    \vec{v}_3 = \vetortri{0}{0}{1}
  \]}
\end{columns}
\end{resolution}

%------------------------------------------------------------------------------%
\begin{resolution}{Resposta}
  \begin{columns}
  \column{0.33\textwidth}
  \[
    \lambda_{1} = 2
  \]
  \[
    \vec{v}_1 = \vetortri{1}{0}{0}
  \]
  \[
    E(2) = \ger\left\{\vec{v}_1\right\}
  \]
  \column{0.33\textwidth}
  \[
    \lambda_{2} = 3
  \]
  \[
    \vec{v}_2 = \vetortri{1}{1}{0}
  \]
  \[
    E(3) = \ger\left\{\vec{v}_2\right\}
  \]
  \column{0.35\textwidth}
  \[
    \lambda_{3} = -1
  \]
  \[
    \vec{v}_3 = \vetortri{0}{0}{1}
  \]
  \[
    E(-1) = \ger\left\{\vec{v}_3\right\}
  \]
  \end{columns}
\end{resolution}

%------------------------------------------------------------------------------%
\endinput

%------------------------------------------------------------------------------%
\begin{question}
  Determine os autovalores, os autovetores e os autoespaços da matriz
  \[
    B =
    \begin{bmatrix}
      4 & 1 & 0 \\
      1 & 4 & 0 \\
      0 & 0 & 2
    \end{bmatrix}
  \]
\end{question}

%------------------------------------------------------------------------------%
\begin{resolution}{Autovalores}
  Polinômio característico
  \begin{align*}
    p(\lambda)
    & = \det(B-\lambda I) \\
    & = \det\left(
          \begin{bmatrix}
            2 & 1 & 0 \\
            0 & 2 & 0 \\
            0 & 0 & 5
          \end{bmatrix}
          -
          \begin{bmatrix}
            \lambda & 0       & 0       \\
            0       & \lambda & 0       \\
            0       & 0       & \lambda
          \end{bmatrix}
          \right)
    \;=\; \begin{vmatrix}
      2-\lambda & 1         & 0         \\
      0         & 2-\lambda & 0         \\
      0         & 0         & 5-\lambda
    \end{vmatrix}
    \\
    & = (2-\lambda)^2(5-\lambda)
  \end{align*}
  Logo, os autovalores são
  \[
    \lambda_{1} = 2
    \qquad
    \lambda_{2} = 5
  \]
\end{resolution}

%------------------------------------------------------------------------------%
\begin{resolution}{Autovetores de \(\lambda=2\)}
\begin{columns}
\column{0.5\textwidth}
  \begin{align*}
    (B-2I)\vec{v} & = \vec{0}
    \\[1mm]
    \begin{bmatrix}[c]
      2-2 & 1   & 0   \\
      0   & 2-2 & 0   \\
      0   & 0   & 5-2
    \end{bmatrix} & =
    \begin{bmatrix}
      0 \\
      0 \\
      0
    \end{bmatrix}
    \\[3mm]
    \begin{bmatrix}
      0 & 1 & 0 \\
      0 & 0 & 0 \\
      0 & 0 & 3
    \end{bmatrix}
    \begin{bmatrix}
      x \\
      y \\
      z
    \end{bmatrix}
    & =
    \begin{bmatrix}
      0 \\
      0 \\
      0
    \end{bmatrix}
   \end{align*}
\column{0.4\textwidth}
  Solução
  \begin{align*}
    x & = t \in \R \\
    y & = 0 \\
    z & = 0
  \end{align*}
  Um autovetor
  \[
    \vec{v}_1
    = \vetortri{1}{0}{0}
  \]
\end{columns}
\end{resolution}

%------------------------------------------------------------------------------%
\begin{resolution}{Autoespaço de \(\lambda=2\)}
  \[
    E(2)
    = \left\{
      \begin{bmatrix}
        t \\
        0 \\
        0
      \end{bmatrix}
      \st t \in \R
    \right\}
    = \ger\left\{
      \begin{bmatrix}
        1 \\
        0 \\
        0
      \end{bmatrix}
    \right\}
  \]
\end{resolution}

%------------------------------------------------------------------------------%
\begin{resolution}{Autovetores de \(\lambda=5\)}
\begin{columns}
\column{0.5\textwidth}
  \begin{align*}
    (B-5I)\vec{v} & = \vec{0}
    \\[1mm]
    \begin{bmatrix}[c]
      2-5 & 1   & 0   \\
      0   & 2-5 & 0   \\
      0   & 0   & 5-5
    \end{bmatrix} & =
    \begin{bmatrix}
      0 \\
      0 \\
      0
    \end{bmatrix}
    \\[3mm]
    \begin{bmatrix}
      -3 & 1  & 0 \\
      0  & -3 & 0 \\
      0  & 0  & 0
    \end{bmatrix}
    \begin{bmatrix}
      x \\
      y \\
      z
    \end{bmatrix}
    & =
    \begin{bmatrix}
      0 \\
      0 \\
      0
    \end{bmatrix}
   \end{align*}
\column{0.4\textwidth}
  Solução
  \begin{align*}
    -3y   & = 0     \\
    -3x+y & = 0     \\
    x     & = y = 0 \\
    z     & = t \in \R
  \end{align*}
  Um autovetor
  \[
    \vec{v}_2
    = \vetortri{0}{0}{1}
  \]
\end{columns}
\end{resolution}

%------------------------------------------------------------------------------%
\begin{resolution}{Autoespaço de \(\lambda=5\)}
  \[
    E_5
    = \left\{
      \begin{bmatrix}
        0 \\
        0 \\
        t
      \end{bmatrix}
      \st t \in \R
    \right\}
    = \ger\left\{
      \begin{bmatrix}
        0 \\
        0 \\
        1
      \end{bmatrix}
    \right\}
  \]
\end{resolution}

%------------------------------------------------------------------------------%
\endinput



\begin{frame}{Exemplo 4: autovalores}
    Os autovalores são as raízes do polinômio característico
    \begin{align*}
        p(\lambda)
        &=\det(B-\lambda I)\\
        &=
        \begin{vmatrix}
            4-\lambda & 1 & 0\\
            1 & 4-\lambda & 0\\
            0 & 0 & 2-\lambda
        \end{vmatrix}\\
        &=(2-\lambda)
        \begin{vmatrix}
            4-\lambda & 1\\
            1 & 4-\lambda
        \end{vmatrix}\\
        &=(2-\lambda)
        \left[(4-\lambda)^2-1\right]\\
        &=(2-\lambda)(3-\lambda)(5-\lambda)
    \end{align*}
\end{frame}

\begin{frame}{Exemplo 4: autovalores}
    Resolvendo
    \[
        (2-\lambda)(3-\lambda)(5-\lambda)=0
    \]

    Obtemos
    \[
        \lambda_{1}=2
        \qquad
        \lambda_{2}=3
        \qquad
        \lambda_{3}=5
    \]
\end{frame}

\begin{frame}{Exemplo 4: autovetores para $\lambda=2$}
    Resolvemos o sistema
    \begin{align*}
        (B-2I)\vec{x}&=\vec{0}\\
        \begin{bmatrix}
            2 & 1 & 0\\
            1 & 2 & 0\\
            0 & 0 & 0
        \end{bmatrix}
        \begin{bmatrix}
            x\\y\\z
        \end{bmatrix}
        &=
        \begin{bmatrix}
            0\\0\\0
        \end{bmatrix}
    \end{align*}

    A forma escalonada do sistema fornece
    \begin{align*}
        2x+y&=0\\
        x+2y&=0\\
        z&\in\mathbb{R}
    \end{align*}

    Logo
    \begin{align*}
        x&=0\\
        y&=0
    \end{align*}

    Portanto
    \[
        E_{2}=
        \operatorname{span}
        \left\{\vetortri{0}{0}{1}\right\}
    \]
\end{frame}

\begin{frame}{Exemplo 4: autovetores para $\lambda=3$}
    Resolvemos o sistema
    \begin{align*}
        (B-3I)\vec{x}&=\vec{0}\\
        \begin{bmatrix}
            1 & 1 & 0\\
            1 & 1 & 0\\
            0 & 0 & -1
        \end{bmatrix}
        \begin{bmatrix}
            x\\y\\z
        \end{bmatrix}
        &=
        \begin{bmatrix}
            0\\0\\0
        \end{bmatrix}
    \end{align*}

    Obtemos
    \begin{align*}
        x+y&=0\\
        z&=0\\
        y&=-x
    \end{align*}

    Portanto
    \[
        E_{3}=
        \left\{
            \vetortri{t}{-t}{0}
            \mid t\in\mathbb{R}
        \right\}
        =
        \operatorname{span}
        \left\{\vetortri{1}{-1}{0}\right\}
    \]
\end{frame}

\begin{frame}{Exemplo 4: autovetores para $\lambda=5$}
    Resolvemos o sistema
    \begin{align*}
        (B-5I)\vec{x}&=\vec{0}\\
        \begin{bmatrix}
            -1 & 1 & 0\\
            1 & -1 & 0\\
            0 & 0 & -3
        \end{bmatrix}
        \begin{bmatrix}
            x\\y\\z
        \end{bmatrix}
        &=
        \begin{bmatrix}
            0\\0\\0
        \end{bmatrix}
    \end{align*}

    Obtemos
    \begin{align*}
        -x+y&=0\\
        z&=0\\
        y&=x
    \end{align*}

    Portanto
    \[
        E_{5}=
        \left\{
            \vetortri{t}{t}{0}
            \mid t\in\mathbb{R}
        \right\}
        =
        \operatorname{span}
        \left\{\vetortri{1}{1}{0}\right\}
    \]
\end{frame}

\begin{frame}{Exemplo 4: resposta}
    \[
        \lambda_{1}=2
        \qquad
        \lambda_{2}=3
        \qquad
        \lambda_{3}=5
    \]

    Os autoespaços são
    \begin{align*}
        E_{2}&=\operatorname{span}
        \left\{\vetortri{0}{0}{1}\right\}\\
        E_{3}&=\operatorname{span}
        \left\{\vetortri{1}{-1}{0}\right\}\\
        E_{5}&=\operatorname{span}
        \left\{\vetortri{1}{1}{0}\right\}
    \end{align*}

    Os autovetores são todos os vetores não nulos pertencentes aos respectivos autoespaços
\end{frame}
%------------------------------------------------------------------------------%
\begin{question}
  Determine os autovalores, os autovetores e os autoespaços da matriz
  \[
    A =
    \begin{pmatrix}
      1 & 1 & 0  \\
      0 & 3 & 0  \\
      0 & 0 & -2
    \end{pmatrix}
  \]
\end{question}

%------------------------------------------------------------------------------%
\begin{resolution}{Autovalores}
  \onslide<+->{Os autovalores são as raízes do polinômio característico}
  \begin{align*}
      \onslide<+->{p(\lambda)}
      \onslide<+->{& = \det(A-\lambda I)}
      \\
      \onslide<+->{& = \begin{vmatrix}
        1-\lambda & 1         & 0          \\
        0         & 3-\lambda & 0          \\
        0         & 0         & -2-\lambda
      \end{vmatrix}}
      \\
      \onslide<+->{& = (1-\lambda)(3-\lambda)(-2-\lambda)}
  \end{align*}

  \onslide<+->{Portanto
  \[
      \lambda_{1} = 1
      \qquad
      \lambda_{2} = 3
      \qquad
      \lambda_{3} = -2
  \]}
\end{resolution}

%------------------------------------------------------------------------------%
\begin{resolution}{Autovetores para $\lambda=1$}
\begin{columns}
\column{0.7\textwidth}
  \begin{align*}
      \onslide<+->{(A - 1I) \vec{x} & = \vec{0}}
      \\[3mm]
      \onslide<+->{\left(
        \begin{bmatrix}
          1 & 1 & 0  \\
          0 & 3 & 0  \\
          0 & 0 & -2
        \end{bmatrix}
      - \begin{bmatrix}
          1 & 0 & 0  \\
          0 & 1 & 0  \\
          0 & 0 & 1
        \end{bmatrix}
      \right)
      \vetortri{x}{y}{z}
      & = \vetortri{0}{0}{0}}
      \\[3mm]
      \onslide<+->{\begin{pmatrix}
        0 & 1 & 0  \\
        0 & 2 & 0  \\
        0 & 0 & -3
      \end{pmatrix}
      \vetortri{x}{y}{z}
      & = \vetortri{0}{0}{0}}
  \end{align*}

\column{0.3\textwidth}
  \begin{align*}
    \onslide<+->{x & \in \R \\}
    \onslide<+->{y & = 0    \\}
    \onslide<+->{z & = 0}
  \end{align*}

  \onslide<+->{Autovetor
  \[
    \vec{v}_1 = \vetortri{1}{0}{0}
  \]}
\end{columns}
\end{resolution}

%------------------------------------------------------------------------------%
\begin{resolution}{Autovetores para $\lambda=3$}
\begin{columns}
\column{0.7\textwidth}
  \begin{align*}
    \onslide<+->{(A - 3I) \vec{x} & = \vec{0}}
    \\[3mm]
    \onslide<+->{\left(
      \begin{bmatrix}
        1 & 1 & 0  \\
        0 & 3 & 0  \\
        0 & 0 & -2
      \end{bmatrix}
    - \begin{bmatrix}
        3 & 0 & 0  \\
        0 & 3 & 0  \\
        0 & 0 & 3
      \end{bmatrix}
    \right)
    \vetortri{x}{y}{z}
    & = \vetortri{0}{0}{0}}
    \\[3mm]
    \onslide<+->{\begin{pmatrix}
      -2 & 1 & 0  \\
      0  & 0 & 0  \\
      0  & 0 & -5
    \end{pmatrix}
    \vetortri{x}{y}{z}
    & = \vetortri{0}{0}{0}}
  \end{align*}

\column{0.3\textwidth}
  \begin{align*}
    \onslide<+->{z     & = 0 \\[4mm]}
    \onslide<+->{-2x+y & = 0 \\}
    \onslide<+->{y     & = 2x  }
  \end{align*}

  \onslide<+->{Autovetor
  \[
    \vec{v}_2 = \vetortri{1}{2}{0}
  \]}
\end{columns}
\end{resolution}

%------------------------------------------------------------------------------%
\begin{resolution}{Autovetores para $\lambda=-2$}
\begin{columns}
\column{0.7\textwidth}
  \begin{align*}
    \onslide<+->{(A-(-2)I)\vec{x}&=\vec{0}}
    \\[3mm]
    \onslide<+->{\left(
      \begin{bmatrix}
        1 & 1 & 0  \\
        0 & 3 & 0  \\
        0 & 0 & -2
      \end{bmatrix}
    + \begin{bmatrix}
        2 & 0 & 0  \\
        0 & 2 & 0  \\
        0 & 0 & 2
      \end{bmatrix}
    \right)
    \vetortri{x}{y}{z}
    & = \vetortri{0}{0}{0}}
    \\[3mm]
    \onslide<+->{\begin{pmatrix}
        3 & 1 & 0\\
        0 & 5 & 0\\
        0 & 0 & 0
    \end{pmatrix}
    \vetortri{x}{y}{z}
    & = \vetortri{0}{0}{0}}
  \end{align*}

\column{0.3\textwidth}
  \begin{align*}
    \onslide<+->{     y & = 0 \\[3mm]}
    \onslide<+->{3x + y & = 0 \\}
    \onslide<+->{     x & = 0 \\[3mm]}
    \onslide<+->{z      & \in\R}
  \end{align*}

  \onslide<+->{Autovetor
  \[
    \vec{v}_3 = \vetortri{0}{0}{1}
  \]}
\end{columns}
\end{resolution}

%------------------------------------------------------------------------------%
\begin{resolution}{Resposta}
  \begin{columns}
  \column{0.33\textwidth}
  \[
    \lambda_{1} = 1
  \]
  \[
    \vec{v}_1 = \vetortri{1}{0}{0}
  \]
  \[
    E(2) = \ger\left\{\vec{v}_1\right\}
  \]
  \column{0.33\textwidth}
  \[
    \lambda_{2} = 3
  \]
  \[
    \vec{v}_2 = \vetortri{1}{2}{0}
  \]
  \[
    E(3) = \ger\left\{\vec{v}_2\right\}
  \]
  \column{0.35\textwidth}
  \[
    \lambda_{3} = -2
  \]
  \[
    \vec{v}_3 = \vetortri{0}{0}{1}
  \]
  \[
    E(-2) = \ger\left\{\vec{v}_3\right\}
  \]
  \end{columns}
\end{resolution}

%------------------------------------------------------------------------------%
\endinput
%%------------------------------------------------------------------------------%
\begin{question}
  Determine os autovalores, os autovetores e os autoespaços da matriz
  \[
    B =
    \begin{bmatrix}
      4 & 1 & 0 \\
      1 & 4 & 0 \\
      0 & 0 & 2
    \end{bmatrix}
  \]
\end{question}

%------------------------------------------------------------------------------%
\begin{resolution}{Autovalores}
  Polinômio característico
  \begin{align*}
    p(\lambda)
    & = \det(B-\lambda I) \\
    & = \det\left(
          \begin{bmatrix}
            4 & 1 & 0 \\
            1 & 4 & 0 \\
            0 & 0 & 2
          \end{bmatrix}
          -
          \begin{bmatrix}
            \lambda & 0       & 0       \\
            0       & \lambda & 0       \\
            0       & 0       & \lambda
          \end{bmatrix}
          \right)
    \;=\; \begin{vmatrix}
      4-\lambda & 1         & 0         \\
      1         & 4-\lambda & 0         \\
      0         & 0         & 2-\lambda
    \end{vmatrix}
    \\
    & = (2-\lambda)
      \begin{vmatrix}
        4-\lambda & 1         \\
        1         & 4-\lambda
      \end{vmatrix}\\
    & = (2 - \lambda) \left[ (4-\lambda)^2 - 1 \right] \\
    & = (2-\lambda)(3-\lambda)(5-\lambda)
  \end{align*}
  Logo, os autovalores são
  \[
    \lambda_{1} = 2
    \qquad
    \lambda_{2} = 5
  \]
\end{resolution}

%------------------------------------------------------------------------------%
\begin{resolution}{Autovetores de \(\lambda=2\)}
\begin{columns}
\column{0.5\textwidth}
  \begin{align*}
    (B-2I)\vec{v} & = \vec{0}
    \\[1mm]
    \begin{bmatrix}[c]
      2-2 & 1   & 0   \\
      0   & 2-2 & 0   \\
      0   & 0   & 5-2
    \end{bmatrix} & =
    \begin{bmatrix}
      0 \\
      0 \\
      0
    \end{bmatrix}
    \\[3mm]
    \begin{bmatrix}
      0 & 1 & 0 \\
      0 & 0 & 0 \\
      0 & 0 & 3
    \end{bmatrix}
    \begin{bmatrix}
      x \\
      y \\
      z
    \end{bmatrix}
    & =
    \begin{bmatrix}
      0 \\
      0 \\
      0
    \end{bmatrix}
   \end{align*}
\column{0.4\textwidth}
  Solução
  \begin{align*}
    x & = t \in \R \\
    y & = 0 \\
    z & = 0
  \end{align*}
  Um autovetor
  \[
    \vec{v}_1
    = \vetortri{1}{0}{0}
  \]
\end{columns}
\end{resolution}

%------------------------------------------------------------------------------%
\begin{resolution}{Autoespaço de \(\lambda=2\)}
  \[
    E(2)
    = \left\{
      \begin{bmatrix}
        t \\
        0 \\
        0
      \end{bmatrix}
      \st t \in \R
    \right\}
    = \ger\left\{
      \begin{bmatrix}
        1 \\
        0 \\
        0
      \end{bmatrix}
    \right\}
  \]
\end{resolution}

%------------------------------------------------------------------------------%
\begin{resolution}{Autovetores de \(\lambda=5\)}
\begin{columns}
\column{0.5\textwidth}
  \begin{align*}
    (B-5I)\vec{v} & = \vec{0}
    \\[1mm]
    \begin{bmatrix}[c]
      2-5 & 1   & 0   \\
      0   & 2-5 & 0   \\
      0   & 0   & 5-5
    \end{bmatrix} & =
    \begin{bmatrix}
      0 \\
      0 \\
      0
    \end{bmatrix}
    \\[3mm]
    \begin{bmatrix}
      -3 & 1  & 0 \\
      0  & -3 & 0 \\
      0  & 0  & 0
    \end{bmatrix}
    \begin{bmatrix}
      x \\
      y \\
      z
    \end{bmatrix}
    & =
    \begin{bmatrix}
      0 \\
      0 \\
      0
    \end{bmatrix}
   \end{align*}
\column{0.4\textwidth}
  Solução
  \begin{align*}
    -3y   & = 0     \\
    -3x+y & = 0     \\
    x     & = y = 0 \\
    z     & = t \in \R
  \end{align*}
  Um autovetor
  \[
    \vec{v}_2
    = \vetortri{0}{0}{1}
  \]
\end{columns}
\end{resolution}

%------------------------------------------------------------------------------%
\begin{resolution}{Autoespaço de \(\lambda=5\)}
  \[
    E_5
    = \left\{
      \begin{bmatrix}
        0 \\
        0 \\
        t
      \end{bmatrix}
      \st t \in \R
    \right\}
    = \ger\left\{
      \begin{bmatrix}
        0 \\
        0 \\
        1
      \end{bmatrix}
    \right\}
  \]
\end{resolution}

%------------------------------------------------------------------------------%
%\endinput

\begin{align*}
---
\end{align*}
%------------------------------------------------------------------------------%
%------------------------------------------------------------------------------%
%------------------------------------------------------------------------------%
%------------------------------------------------------------------------------%
%------------------------------------------------------------------------------%


\begin{frame}{Exemplo 4: autovalores}
    Os autovalores são as raízes do polinômio característico
    \begin{align*}
        p(\lambda)
        &=\det(B-\lambda I)\\
        &=
        \begin{vmatrix}
            4-\lambda & 1 & 0\\
            1 & 4-\lambda & 0\\
            0 & 0 & 2-\lambda
        \end{vmatrix}\\
        &=(2-\lambda)
        \begin{vmatrix}
            4-\lambda & 1\\
            1 & 4-\lambda
        \end{vmatrix}\\
        &=(2-\lambda)
        \left[(4-\lambda)^2-1\right]\\
        &=(2-\lambda)(3-\lambda)(5-\lambda)
    \end{align*}
\end{frame}

\begin{frame}{Exemplo 4: autovalores}
    Resolvendo
    \[
        (2-\lambda)(3-\lambda)(5-\lambda)=0
    \]

    Obtemos
    \[
        \lambda_{1}=2
        \qquad
        \lambda_{2}=3
        \qquad
        \lambda_{3}=5
    \]
\end{frame}

\begin{frame}{Exemplo 4: autovetores para $\lambda=2$}
    Resolvemos o sistema
    \begin{align*}
        (B-2I)\vec{x}&=\vec{0}\\
        \begin{bmatrix}
            2 & 1 & 0\\
            1 & 2 & 0\\
            0 & 0 & 0
        \end{bmatrix}
        \begin{bmatrix}
            x\\y\\z
        \end{bmatrix}
        &=
        \begin{bmatrix}
            0\\0\\0
        \end{bmatrix}
    \end{align*}

    A forma escalonada do sistema fornece
    \begin{align*}
        2x+y&=0\\
        x+2y&=0\\
        z&\in\mathbb{R}
    \end{align*}

    Logo
    \begin{align*}
        x&=0\\
        y&=0
    \end{align*}

    Portanto
    \[
        E_{2}=
        \operatorname{span}
        \left\{\vetortri{0}{0}{1}\right\}
    \]
\end{frame}

\begin{frame}{Exemplo 4: autovetores para $\lambda=3$}
    Resolvemos o sistema
    \begin{align*}
        (B-3I)\vec{x}&=\vec{0}\\
        \begin{bmatrix}
            1 & 1 & 0\\
            1 & 1 & 0\\
            0 & 0 & -1
        \end{bmatrix}
        \begin{bmatrix}
            x\\y\\z
        \end{bmatrix}
        &=
        \begin{bmatrix}
            0\\0\\0
        \end{bmatrix}
    \end{align*}

    Obtemos
    \begin{align*}
        x+y&=0\\
        z&=0\\
        y&=-x
    \end{align*}

    Portanto
    \[
        E_{3}=
        \left\{
            \vetortri{t}{-t}{0}
            \mid t\in\mathbb{R}
        \right\}
        =
        \operatorname{span}
        \left\{\vetortri{1}{-1}{0}\right\}
    \]
\end{frame}

\begin{frame}{Exemplo 4: autovetores para $\lambda=5$}
    Resolvemos o sistema
    \begin{align*}
        (B-5I)\vec{x}&=\vec{0}\\
        \begin{bmatrix}
            -1 & 1 & 0\\
            1 & -1 & 0\\
            0 & 0 & -3
        \end{bmatrix}
        \begin{bmatrix}
            x\\y\\z
        \end{bmatrix}
        &=
        \begin{bmatrix}
            0\\0\\0
        \end{bmatrix}
    \end{align*}

    Obtemos
    \begin{align*}
        -x+y&=0\\
        z&=0\\
        y&=x
    \end{align*}

    Portanto
    \[
        E_{5}=
        \left\{
            \vetortri{t}{t}{0}
            \mid t\in\mathbb{R}
        \right\}
        =
        \operatorname{span}
        \left\{\vetortri{1}{1}{0}\right\}
    \]
\end{frame}

\begin{frame}{Exemplo 4: resposta}
    \[
        \lambda_{1}=2
        \qquad
        \lambda_{2}=3
        \qquad
        \lambda_{3}=5
    \]

    Os autoespaços são
    \begin{align*}
        E_{2}&=\operatorname{span}
        \left\{\vetortri{0}{0}{1}\right\}\\
        E_{3}&=\operatorname{span}
        \left\{\vetortri{1}{-1}{0}\right\}\\
        E_{5}&=\operatorname{span}
        \left\{\vetortri{1}{1}{0}\right\}
    \end{align*}

    Os autovetores são todos os vetores não nulos pertencentes aos respectivos autoespaços
\end{frame}
%\begin{question}
    Determine os autovalores, os autovetores e os autoespaços da matriz
    \[
        B=
        \begin{pmatrix}
            4 & 1 & 1\\
            1 & 4 & 1\\
            1 & 1 & 4
        \end{pmatrix}
    \]
\end{question}

\begin{resolution}{Exemplo 4: autovalores}
    O polinômio característico é
    \begin{align*}
        p(\lambda)
        &=\det(B-\lambda I)\\
        &=
        \begin{vmatrix}
            4-\lambda & 1 & 1\\
            1 & 4-\lambda & 1\\
            1 & 1 & 4-\lambda
        \end{vmatrix}\\
        &=(4-\lambda)^3-3(4-\lambda)+2\\
        &=(6-\lambda)(3-\lambda)^2
    \end{align*}

    Portanto
    \[
        \lambda_{1}=6
        \qquad
        \lambda_{2}=3
    \]
\end{resolution}

\begin{resolution}{Exemplo 4: autovetores para $\lambda=6$}
    Resolvemos o sistema
    \begin{align*}
        (B-6I)\vec{x}&=\vec{0}\\
        \begin{pmatrix}
            -2 & 1 & 1\\
            1 & -2 & 1\\
            1 & 1 & -2
        \end{pmatrix}
        \begin{pmatrix}
            x\\y\\z
        \end{pmatrix}
        &=
        \begin{pmatrix}
            0\\0\\0
        \end{pmatrix}
    \end{align*}

    O sistema equivalente é
    \begin{align*}
        -2x+y+z&=0\\
        x-2y+z&=0\\
        x+y-2z&=0
    \end{align*}

    Subtraindo as equações, obtemos
    \begin{align*}
        x&=y\\
        y&=z
    \end{align*}
\end{resolution}

\begin{resolution}{Exemplo 4: autoespaço para $\lambda=6$}
    Portanto
    \[
        E_{6}=
        \left\{
            \vetortri{t}{t}{t}
            \mid t\in\mathbb{R}
        \right\}
    \]

    Ou seja
    \[
        E_{6}=
        \operatorname{span}
        \left\{\vetortri{1}{1}{1}\right\}
    \]

    Os autovetores são os múltiplos não nulos de $\vetortri{1}{1}{1}$
\end{resolution}

\begin{resolution}{Exemplo 4: autovetores para $\lambda=3$}
    Resolvemos o sistema
    \begin{align*}
        (B-3I)\vec{x}&=\vec{0}\\
        \begin{pmatrix}
            1 & 1 & 1\\
            1 & 1 & 1\\
            1 & 1 & 1
        \end{pmatrix}
        \begin{pmatrix}
            x\\y\\z
        \end{pmatrix}
        &=
        \begin{pmatrix}
            0\\0\\0
        \end{pmatrix}
    \end{align*}

    O sistema equivalente é
    \[
        x+y+z=0
    \]

    Escolhendo $x=s$ e $y=t$, temos
    \[
        z=-s-t
    \]
\end{resolution}

\begin{resolution}{Exemplo 4: autoespaço para $\lambda=3$}
    Portanto
    \[
        E_{3}=
        \left\{
            \vetortri{s}{t}{-s-t}
            \mid s,t\in\mathbb{R}
        \right\}
    \]

    Escrevendo como combinação linear
    \begin{align*}
        E_{3}
        &=
        \left\{
            s\vetortri{1}{0}{-1}
            +
            t\vetortri{0}{1}{-1}
            \mid s,t\in\mathbb{R}
        \right\}\\
        &=
        \operatorname{span}
        \left\{
            \vetortri{1}{0}{-1}
            ,
            \vetortri{0}{1}{-1}
        \right\}
    \end{align*}

    Logo
    \[
        \dim(E_{3})=2
    \]
\end{resolution}

\begin{resolution}{Exemplo 4: resposta}
    Os autovalores são
    \[
        \lambda_{1}=6
        \qquad
        \lambda_{2}=3
    \]

    Os autoespaços são
    \begin{align*}
        E_{6}&=
        \operatorname{span}
        \left\{\vetortri{1}{1}{1}\right\}\\
        E_{3}&=
        \operatorname{span}
        \left\{
            \vetortri{1}{0}{-1}
            ,
            \vetortri{0}{1}{-1}
        \right\}
    \end{align*}

    O autoespaço $E_{3}$ possui dimensão dois
\end{resolution}

%%------------------------------------------------------------------------------%
%\section{Lista Mínima}
%
%%------------------------------------------------------------------------------%
%\begin{frame}{Lista Mínima}
%  \begin{enumerate}
%    \item
%  \end{enumerate}
%
%  \vfill
%  Atenção: A prova é baseada no livro, não nas apresentações
%\end{frame}

%------------------------------------------------------------------------------%
\end{document}
%------------------------------------------------------------------------------%
